% Options for packages loaded elsewhere
% Options for packages loaded elsewhere
\PassOptionsToPackage{unicode}{hyperref}
\PassOptionsToPackage{hyphens}{url}
\PassOptionsToPackage{dvipsnames,svgnames,x11names}{xcolor}
%
\documentclass[
  11pt,
]{article}
\usepackage{xcolor}
\usepackage[margin=1in]{geometry}
\usepackage{amsmath,amssymb}
\setcounter{secnumdepth}{-\maxdimen} % remove section numbering
\usepackage{iftex}
\ifPDFTeX
  \usepackage[T1]{fontenc}
  \usepackage[utf8]{inputenc}
  \usepackage{textcomp} % provide euro and other symbols
\else % if luatex or xetex
  \usepackage{unicode-math} % this also loads fontspec
  \defaultfontfeatures{Scale=MatchLowercase}
  \defaultfontfeatures[\rmfamily]{Ligatures=TeX,Scale=1}
\fi
\usepackage{lmodern}
\ifPDFTeX\else
  % xetex/luatex font selection
\fi
% Use upquote if available, for straight quotes in verbatim environments
\IfFileExists{upquote.sty}{\usepackage{upquote}}{}
\IfFileExists{microtype.sty}{% use microtype if available
  \usepackage[]{microtype}
  \UseMicrotypeSet[protrusion]{basicmath} % disable protrusion for tt fonts
}{}
\makeatletter
\@ifundefined{KOMAClassName}{% if non-KOMA class
  \IfFileExists{parskip.sty}{%
    \usepackage{parskip}
  }{% else
    \setlength{\parindent}{0pt}
    \setlength{\parskip}{6pt plus 2pt minus 1pt}}
}{% if KOMA class
  \KOMAoptions{parskip=half}}
\makeatother
% Make \paragraph and \subparagraph free-standing
\makeatletter
\ifx\paragraph\undefined\else
  \let\oldparagraph\paragraph
  \renewcommand{\paragraph}{
    \@ifstar
      \xxxParagraphStar
      \xxxParagraphNoStar
  }
  \newcommand{\xxxParagraphStar}[1]{\oldparagraph*{#1}\mbox{}}
  \newcommand{\xxxParagraphNoStar}[1]{\oldparagraph{#1}\mbox{}}
\fi
\ifx\subparagraph\undefined\else
  \let\oldsubparagraph\subparagraph
  \renewcommand{\subparagraph}{
    \@ifstar
      \xxxSubParagraphStar
      \xxxSubParagraphNoStar
  }
  \newcommand{\xxxSubParagraphStar}[1]{\oldsubparagraph*{#1}\mbox{}}
  \newcommand{\xxxSubParagraphNoStar}[1]{\oldsubparagraph{#1}\mbox{}}
\fi
\makeatother


\usepackage{longtable,booktabs,array}
\newcounter{none} % for unnumbered tables
\usepackage{calc} % for calculating minipage widths
% Correct order of tables after \paragraph or \subparagraph
\usepackage{etoolbox}
\makeatletter
\patchcmd\longtable{\par}{\if@noskipsec\mbox{}\fi\par}{}{}
\makeatother
% Allow footnotes in longtable head/foot
\IfFileExists{footnotehyper.sty}{\usepackage{footnotehyper}}{\usepackage{footnote}}
\makesavenoteenv{longtable}
\usepackage{graphicx}
\makeatletter
\newsavebox\pandoc@box
\newcommand*\pandocbounded[1]{% scales image to fit in text height/width
  \sbox\pandoc@box{#1}%
  \Gscale@div\@tempa{\textheight}{\dimexpr\ht\pandoc@box+\dp\pandoc@box\relax}%
  \Gscale@div\@tempb{\linewidth}{\wd\pandoc@box}%
  \ifdim\@tempb\p@<\@tempa\p@\let\@tempa\@tempb\fi% select the smaller of both
  \ifdim\@tempa\p@<\p@\scalebox{\@tempa}{\usebox\pandoc@box}%
  \else\usebox{\pandoc@box}%
  \fi%
}
% Set default figure placement to htbp
\def\fps@figure{htbp}
\makeatother





\setlength{\emergencystretch}{3em} % prevent overfull lines

\providecommand{\tightlist}{%
  \setlength{\itemsep}{0pt}\setlength{\parskip}{0pt}}



 


\usepackage{amsmath,amssymb,enumitem,booktabs}
\setlist{itemsep=2pt,topsep=3pt}
\usepackage{fancyhdr}\pagestyle{fancy}\fancyhf{}
\fancyhead[L]{\small DS701 Fall 2026}\fancyhead[R]{\small Homework 10}
\fancyfoot[C]{\small\thepage}
\renewcommand{\headrulewidth}{0.4pt}
\usepackage[breakable]{tcolorbox}
% --- answer-space macros (problems PDF only; no-ops in the solutions PDF) ---
% \answerspace{2in}  vertical room for handwritten work, sized from the solution
% \blank             a short fill-in rule; \blank[1.5in] for a wider one
% \answerpage        start the next problem on a fresh page (Gradescope page matching)
\newcommand{\answerspace}[1]{\par\vspace*{#1}}  % \vspace* so room is not dropped at a page break
\newcommand{\blank}[1][1.2in]{\rule{#1}{0.4pt}}
\newcommand{\answerpage}{\clearpage}
\newtcolorbox{solutionbox}{breakable,colback=blue!4,colframe=blue!45!black,title=Solution,fonttitle=\bfseries,left=4pt,right=4pt,top=3pt,bottom=3pt}
\makeatletter
\@ifpackageloaded{caption}{}{\usepackage{caption}}
\AtBeginDocument{%
\ifdefined\contentsname
  \renewcommand*\contentsname{Table of contents}
\else
  \newcommand\contentsname{Table of contents}
\fi
\ifdefined\listfigurename
  \renewcommand*\listfigurename{List of Figures}
\else
  \newcommand\listfigurename{List of Figures}
\fi
\ifdefined\listtablename
  \renewcommand*\listtablename{List of Tables}
\else
  \newcommand\listtablename{List of Tables}
\fi
\ifdefined\figurename
  \renewcommand*\figurename{Figure}
\else
  \newcommand\figurename{Figure}
\fi
\ifdefined\tablename
  \renewcommand*\tablename{Table}
\else
  \newcommand\tablename{Table}
\fi
}
\@ifpackageloaded{float}{}{\usepackage{float}}
\floatstyle{ruled}
\@ifundefined{c@chapter}{\newfloat{codelisting}{h}{lop}}{\newfloat{codelisting}{h}{lop}[chapter]}
\floatname{codelisting}{Listing}
\newcommand*\listoflistings{\listof{codelisting}{List of Listings}}
\makeatother
\makeatletter
\makeatother
\makeatletter
\@ifpackageloaded{caption}{}{\usepackage{caption}}
\@ifpackageloaded{subcaption}{}{\usepackage{subcaption}}
\makeatother
\usepackage{bookmark}
\IfFileExists{xurl.sty}{\usepackage{xurl}}{} % add URL line breaks if available
\urlstyle{same}
% fallback for those not using the hyperref driver hyperxmp:
\makeatletter
\@ifundefined{xmpquote}{\newcommand{\xmpquote}[1]{#1}}{}
\makeatother
\hypersetup{
  pdftitle={DS701 Homework 10: TF-IDF{,} Matrix Factorization{,} and Graph Basics},
  colorlinks=true,
  linkcolor={blue},
  filecolor={Maroon},
  citecolor={Blue},
  urlcolor={Blue},
  pdfcreator={LaTeX via pandoc}}


\title{DS701 Homework 10: TF-IDF, Matrix Factorization, and Graph
Basics}
\usepackage{etoolbox}
\makeatletter
\providecommand{\subtitle}[1]{% add subtitle to \maketitle
  \apptocmd{\@title}{\par {\large #1 \par}}{}{}
}
\makeatother
\subtitle{Released Mon Nov 16 · due Mon Nov 23, 11:59 pm on Gradescope}
\author{}
\date{}
\begin{document}
\maketitle


\noindent\textbf{Name:}\ \rule{2.0in}{0.4pt}\qquad\textbf{BU ID:}\ \rule{1.2in}{0.4pt}
\vspace{6pt}

\textbf{How this works.} Answer \textbf{by hand on paper} (or a tablet
with a stylus), show your work, then scan and upload to Gradescope,
matching each problem to its pages. This is deliberate practice for the
closed-notes quizzes, which are in the same style --- so do it as much
as possible without a computer or an AI assistant. If you are stuck, you
can try AI in tutorial style or reach out to course staff. It is
\textbf{participation-graded}: a serious attempt at every problem earns
full credit, and worked solutions are released the day after the
deadline. Three problems, two parts each; written for well under two
hours.

\textbf{Covers:} Lecture 18 (NLP: text to vectors, Mon Nov 16), Lecture
19 (Recommender Systems, Wed Nov 18) and Lecture 20 (Networks I, Mon Nov
23 --- its recording is posted ahead of the session, so watch it before
the due date). Anything you need beyond the lectures (a table, a model's
output, a precomputed weight) is printed in the problem. Where
arithmetic is asked for, it is one step; show that one line of work.

\subsection{Problem 1: Reading a TF-IDF
table}\label{problem-1-reading-a-tf-idf-table}

\emph{Practices: what a TF-IDF weight says, why common words vanish, one
cosine similarity, and what a dense embedding adds (L18).}

A corpus holds \(N = 8\) customer reviews of a phone, lower-cased and
tokenized on spaces. The table shows, for two of the reviews and four
vocabulary words, the term count \(\text{tf}(t,d)\), the document
frequency \(\text{df}(t)\) (how many of the 8 reviews contain the word),
the resulting \(\text{idf}(t) = \log_2\frac{N}{\text{df}(t)}\) (base 2,
so the values are whole numbers), and the TF-IDF weight
\(\text{tf}\times\text{idf}\).

\begin{center}
\begin{tabular}{@{}l cccc@{}}\toprule
 & \texttt{battery} & \texttt{screen} & \texttt{the} & \texttt{refund}\\ \midrule
tf in review 1 & 2 & 3 & 5 & 0\\
tf in review 2 & 0 & 4 & 4 & 1\\ \midrule
df (out of $N=8$) & 2 & 4 & 8 & 1\\
idf $=\log_2(8/\text{df})$ & 2 & 1 & 0 & 3\\ \midrule
TF-IDF, review 1 & 4 & 3 & 0 & 0\\
TF-IDF, review 2 & 0 & 4 & 0 & 3\\
\bottomrule\end{tabular}\end{center}

\begin{enumerate}
\def\labelenumi{(\alph{enumi})}
\tightlist
\item
  \textbf{Read the table.} (i) \texttt{the} is the most frequent word in
  both reviews, yet its TF-IDF weight is 0 in both. Say in one sentence
  why, and why that is the intended behaviour. (ii) In review 2,
  \texttt{refund} appears only \emph{once} yet earns a weight of 3,
  nearly as much as \texttt{screen} gets from four occurrences ---
  explain in a sentence what a large TF-IDF weight means (``frequent
  here, rare across the corpus''). (iii) A third review says ``cracked
  display'' and never uses the word \texttt{screen}. Would a TF-IDF
  search for the query \texttt{screen} find it? What would a dense word
  embedding (Word2Vec, GloVe) let the system see that the sparse vector
  cannot? One or two sentences.
\end{enumerate}

\begin{enumerate}
\def\labelenumi{(\roman{enumi})}
\item
  \answerspace{0.8in}
\item
  \answerspace{0.8in}
\item
  \answerspace{1.0in}
\end{enumerate}

\begin{enumerate}
\def\labelenumi{(\alph{enumi})}
\setcounter{enumi}{1}
\tightlist
\item
  \textbf{One cosine.} Using the two TF-IDF rows as 4-vectors, compute
  the cosine similarity between review 1 and review 2 (one line of work;
  a fraction is fine). Then, \emph{without computing it}, say whether
  the cosine of the \textbf{raw count} rows would come out higher or
  lower, and which word is responsible.
\end{enumerate}

\(\cos(\text{review 1}, \text{review 2})\) using TF-IDF =
\blank[0.9in]\qquad raw counts would be higher / lower (circle one),
because of \blank[1.0in]

\answerspace{1.6in}

\answerpage

\subsection{Problem 2: A ratings matrix and a matrix-factorization
model}\label{problem-2-a-ratings-matrix-and-a-matrix-factorization-model}

\emph{Practices: which entries the MF objective sums over, reading a
held-out RMSE table, one factor dot product, and cold start (L19).}

Four users have rated some of four films on a 1--5 scale; a dash means
\emph{no rating} (left). A rank-\(k\) matrix-factorization model
predicts \(\hat r_{ui} = u_u^T v_i\) and is fit by minimising
\(\sum (r_{ui} - u_u^Tv_i)^2 + \lambda(\|U\|^2 + \|V\|^2)\); it was
trained for several \(k\) with the ratings split into a training and a
held-out set, with the RMSEs on the right.

\begin{center}
\begin{minipage}[c]{0.42\textwidth}\centering
\begin{tabular}{@{}l cccc@{}}\toprule
 & $F_1$ & $F_2$ & $F_3$ & $F_4$\\ \midrule
Ana & 5 & -- & 4 & --\\
Ben & -- & 3 & -- & 2\\
Cal & 4 & 1 & -- & --\\
Dee & -- & -- & 5 & 3\\
\bottomrule\end{tabular}
\end{minipage}\hfill
\begin{minipage}[c]{0.54\textwidth}\centering
\begin{tabular}{@{}l cccc@{}}\toprule
$k$ & 1 & 2 & 4 & 8\\ \midrule
training RMSE & 0.95 & 0.80 & 0.62 & 0.35\\
held-out RMSE & 0.98 & 0.86 & 0.84 & 0.97\\
\bottomrule\end{tabular}
\end{minipage}
\end{center}

\begin{enumerate}
\def\labelenumi{(\alph{enumi})}
\tightlist
\item
  \textbf{Read the output.} (i) How many terms does the squared-error
  sum in the objective have for \emph{this} matrix, and which entries
  are they? Why can the sum not run over all 16 entries --- and what
  does that have to do with why plain SVD cannot be applied to \(R\)
  directly? (ii) From the RMSE table, which \(k\) would you choose, and
  what is the pattern at \(k=8\) called? One sentence each.
\end{enumerate}

\begin{enumerate}
\def\labelenumi{(\roman{enumi})}
\tightlist
\item
  number of terms: \blank[0.7in]
\end{enumerate}

\answerspace{1.0in}

\begin{enumerate}
\def\labelenumi{(\roman{enumi})}
\setcounter{enumi}{1}
\tightlist
\item
  choose \(k\) = \blank[0.5in]\qquad the pattern at \(k=8\) is called
  \blank[1.2in]
\end{enumerate}

\answerspace{0.4in}

\begin{enumerate}
\def\labelenumi{(\alph{enumi})}
\setcounter{enumi}{1}
\tightlist
\item
  \textbf{One prediction, and a new user.} The fitted \(k=2\) factors
  include \(u_{\text{Ben}} = (0.8,\ 1.2)\) and
  \(v_{F_3} = (2.5,\ 1.0)\). (i) What rating does the model predict for
  Ben on \(F_3\) (one multiply-add)? (ii) A new user, Eve, signs up and
  has rated nothing. What does the fitted objective make her factor
  \(u_{\text{Eve}}\), what does the model then predict for her on every
  film, and what is this problem called? Name one way the lecture
  suggested getting around it.
\end{enumerate}

\begin{enumerate}
\def\labelenumi{(\roman{enumi})}
\tightlist
\item
  \(\hat r_{\text{Ben},F_3}\) = \blank[0.9in]
\end{enumerate}

\answerspace{0.3in}

\begin{enumerate}
\def\labelenumi{(\roman{enumi})}
\setcounter{enumi}{1}
\tightlist
\item
  \(u_{\text{Eve}}\) = \blank[0.9in]\qquad predicts for every film:
  \blank[1.6in]\qquad called \blank[1.0in]
\end{enumerate}

\answerspace{0.6in}

\answerpage

\subsection{Problem 3: Graph basics from an edge
list}\label{problem-3-graph-basics-from-an-edge-list}

\emph{Practices: degree, the handshake count, density, one row of the
adjacency matrix, connectivity, and what an entry of \(A^2\) counts
(L20).}

Five buildings on a campus are joined by fibre links. As an undirected
graph on nodes \(\{1,2,3,4,5\}\) the edge set is

\[
E = \{(1,2),\ (1,3),\ (2,3),\ (3,4),\ (4,5)\}.
\]

\begin{enumerate}
\def\labelenumi{(\alph{enumi})}
\tightlist
\item
  \textbf{Small numbers.} Sketch the graph. Give the degree of every
  node, check that the degrees sum to \(2|E|\), compute the density
  \(|E|/\binom{5}{2}\), and write \textbf{row 3} of the adjacency matrix
  \(A\) (columns in order \(1,\dots,5\)).
\end{enumerate}

\(\deg(1)\) = \blank[0.5in]\quad \(\deg(2)\) =
\blank[0.5in]\quad \(\deg(3)\) = \blank[0.5in]\quad \(\deg(4)\) =
\blank[0.5in]\quad \(\deg(5)\) = \blank[0.5in]\textbackslash{[}4pt{]}
\(\sum_i \deg(i)\) = \blank[0.6in]\qquad \(2|E|\) =
\blank[0.6in]\qquad density = \blank[0.8in]\textbackslash{[}4pt{]} row 3
of \(A\):
\(\big(\ \blank[0.3in]\ ,\ \blank[0.3in]\ ,\ \blank[0.3in]\ ,\ \blank[0.3in]\ ,\ \blank[0.3in]\ \big)\)

\answerspace{2.2in}

\begin{enumerate}
\def\labelenumi{(\alph{enumi})}
\setcounter{enumi}{1}
\tightlist
\item
  \textbf{Reading, not multiplying.} (i) Is the graph connected? Name
  one edge whose removal would disconnect it. (ii) Without multiplying
  any matrices: what does the entry \((A^2)_{ij}\) count? Use that to
  read off \((A^2)_{1,5}\) and \((A^2)_{1,4}\) from your sketch (name
  the walk if there is one), and say what \((A^2)_{3,3}\) must equal and
  why.
\end{enumerate}

\begin{enumerate}
\def\labelenumi{(\roman{enumi})}
\tightlist
\item
  connected? \blank[0.6in]\qquad an edge whose removal disconnects it:
  \blank[0.8in]
\end{enumerate}

\answerspace{0.5in}

\begin{enumerate}
\def\labelenumi{(\roman{enumi})}
\setcounter{enumi}{1}
\tightlist
\item
  \((A^2)_{ij}\) counts \blank[2.4in]\textbackslash{[}4pt{]}
  \((A^2)_{1,5}\) = \blank[0.5in]\qquad \((A^2)_{1,4}\) =
  \blank[0.5in]\qquad \((A^2)_{3,3}\) = \blank[0.5in]
\end{enumerate}

\answerspace{1.4in}




\end{document}
